\section{Preliminary results}

\label{sec:results}

Prior to 2022, seating was arranged on a first-come, first-served basis at the Palacio Nacional, which led to criticism as many journalists felt that the process was biased, favoring those aligned with the government over independent or opposition voices\citep{Tombola-Razon}. These journalists would often arrive as early as 1 a.m. to secure front-row seats. In response, the President proposed a lottery system to determine seating. By April 2022, the government formalized this change, announcing that journalists would be randomly assigned to the first three rows, with one journalist per row from each type of media—radio, television, print, etc \citep{Tombola-Financiero}. 

An important assumption in using the lottery assignment to instrument the number of critical or favorable articles written by news outlets is demonstrating that the lottery assignment actually influences which outlets get to ask questions during the press conference. If the president continues to select media outlets that are not critical of the government after the lottery—outlets that also receive more advertising revenue—two scenarios could emerge.

First, the first stage and the reduced form could become insignificant because the lottery assignment would not affect the behavior of journalists. If journalists perceive they have no real chance of participating, their seating row would not alter their approach.


\subsection{Importance of Row Seating}

To examine this, I divided the analysis into two periods: one year before the implementation of the lottery and one year after its implementation. I randomly selected 10 press conferences (five from each period) and recorded data on which journalists asked questions and where they were seated. Table \ref{tab:lottery} presents the descriptive statistics for both periods, offering insights into the changes, if any, in seating and question-asking dynamics. First thing to emphasize is that first row seating is crucial for asking questions in both periods of time. Around 70\% of the questions asked come from journalists seated in the front row. Also, this initial analysis suggests that being seated in the the third or fourth row basically decreases your chances of asking a question by a lot, which could also elicit some behavior on journalists.

\begin{table}[!htpb]
    \centering
    \caption{Descriptive Statistics of Questions Asked}
    \label{tab:lottery}
    \begin{tabular}{cccccc}
\hline \\
Time Period & Total Questions & First Row & Second Row & Third Row & Fourth Row \\
\hline\\
 Before lottery & 27 & 18 (66.67\%) & 7 (25.92\%) & 2 (7.42\%) & 0 \\
 After Lottery & 18 & 14 (77.78\%) & 2 (11.11\%) & 0 & 2 (11.11\%) \\
 Both & 45 & 32 (71.11\%) & 9 (20\%) & 2 (4.44\%) & 2 (4.44\%) \\
\hline \\[-1.8ex]  \multicolumn{6}{l} {\parbox[t]{15cm}{\small \textit{Notes:}  Descriptive statistics of the morning press conferences revised at random from the two relevant time periods. Before the lottery refers to 1 year before April 2022, and after the lottery refers to one year after. Total questions is the total number of questions asked on the respective press conference for each time period. Columns 3-6 refer to how many questions were asked from each respective row.}} \\
\end{tabular}
\end{table}

It is important to note that I am not assuming the two periods are directly comparable, as they are inherently different. For example, based on my observations of the press conferences, the dynamics in 2021 were significantly different from those in 2023. In 2021, COVID-19 safety measures, such as social distancing, face masks, and limits on the number of attendees in enclosed spaces, were in effect. These restrictions led to noticeable differences, such as a significantly smaller number of journalists present in 2021 compared to 2023, when these measures were no longer in place. This for example, could explain why there was a decrease in the total number of questions asked across periods.

Despite these changes, Table \ref{tab:lottery} highlights one aspect that remained consistent across periods: the preference given to first-row journalists when asking questions. Additionally, it is important to emphasize that my analysis will only focus on the period after the lottery was implemented. The pre-lottery period is included only for initial comparisons of select variables that may help strengthen the validity of the instrument.

